\SetRegion{第二章 · 全球开源生态观测地图}{blue}{globe}
\setcounter{chapter}{2}\setcounter{figure}{0}\setcounter{table}{0}
\renewcommand{\thefigure}{2-\arabic{figure}}\renewcommand{\thetable}{2-\arabic{table}}
\phantomsection\label{ch2-home}\pdfbookmark[1]{第二章阅读首页}{ch2-home-bm}
\kicker{02 / GLOBAL OPEN SOURCE OBSERVATORY}
\PageTitle{参与在扩大，\\生态并不同步}{从空间分布，看见规模、结构和增长的差异}
第一章说明了从开放协作走向公共价值需要哪些条件。本章转向可观察的参与与资源：谁进入生态，贡献如何汇聚，项目与人才在哪里积累，以及共同成果怎样获得持续支持。

只看人数、贡献、项目影响或资产供给，可能得到不同的排序和结构。\key{七个主题共同理解一个生态，各自回答明确的问题。}
\par\vspace{4pt}\noindent\includegraphics[width=\linewidth]{figures/new/chapter2_navigation.pdf}
\vfill
\src{主统计期：2025-07-01至2026-06-30。自然年、外部年度窗与独立快照随图注明；时间不明的资料不自动归入主统计期。}
\clearpage
\phantomsection\label{ch2-reader}\pdfbookmark[1]{阅读说明：对象、图形与比较范围}{ch2-reader-bm}
\kicker{READING KEY / 阅读说明}
\PageTitle{先分清对象，\\再比较数字}{一张地图定位，一个读数对应一个范围}
\noindent\includegraphics[width=\linewidth]{figures/new/reading_objects.pdf}
\topic{地图、条形与流带}
地图帮助定位，条形便于比较，趋势图显示变化。流带宽度表示OpenRank贡献度，\key{不表示人员迁移或资金流动}；未知、未列与零分别处理。份额从10\%到12\%，增加的是2个百分点，相对增幅为20\%。
\topic{五域与中国区域}
五域固定比较欧洲、美国、中国、印度与南方六国代表组；后者为巴西、印度尼西亚、越南、菲律宾、埃及和南非。欧洲沿用各表范围，五域不是世界的完整划分。中国省级观察纳入台湾、香港和澳门；城市样本范围另行说明。
\vfill
\src{OpenDigger提供数据与指标基础，OpenShare承载数据访问与应用。OpenRank定义与聚合方式参见\href{https://open-digger.cn/docs/user-docs/metrics/community-openrank}{社区贡献度}与\href{https://open-digger.cn/docs/user-docs/metrics/global-openrank}{全域影响力}。图下来源编号可点击，完整清单与底表随源码保存。}
\clearpage
\phantomsection\label{ch2-windows}\pdfbookmark[1]{时间窗口与样本索引}{ch2-windows-bm}
\kicker{SCOPE / 时间与样本}
\PageTitle{在各自范围内，\\理解比较}{主统计期固定，例外窗口清楚标注}
本报告默认采用\key{2025-07-01至2026-06-30}。图中的“2026报告年度”指该窗口；采集日期、自然年快照和其他平台的年度窗分别记录。

\begingroup\sffamily\fontsize{10.3}{15.4}\selectfont\setlength{\tabcolsep}{4pt}\renewcommand{\arraystretch}{1.45}
\noindent\begin{tabularx}{\linewidth}{@{}p{26mm}X p{53mm}@{}}\toprule
观察内容 & 实际窗口与对象 & 阅读位置\\\midrule
参与与活跃 & 2022—2026报告年度；GitHub地理表、独立五域汇总、本土补充分开。 & 2.1；三个中国活跃读数不能互换。\\\addlinespace[5pt]
贡献与项目 & 全球及省级主要为2025自然年；五域及贡献来源为报告年度。 & 2.2、2.3；不跨表计算转化率。\\\addlinespace[5pt]
人才 & 全球分布为2026年7月快照；五域与固定十城为报告年度序列。 & 2.4；前5\%筛选细节待核。\\\addlinespace[5pt]
代码与HF样本 & 代码为报告年度活跃仓库；发布单位为日期待核的累计样本；任务快照为2026-10-15。 & 2.5；下载不等于独立用户。\\\addlinespace[5pt]
外部平台观察 & GitHub年度活动窗为2025-10至2026-09；语言同比为2025-09与2026-09。HF夏季观察为2026-01至08，具体分析另注。 & 2.5、2.6；不并入本报告主期。\\\addlinespace[5pt]
治理与许可 & 2022—2026报告年度归档；协作材料来自GitHub 2025报告；规范文本核对日期为2026-11-05。 & 2.7；组织、文件与规则不直接证明治理成效。\\\bottomrule
\end{tabularx}\endgroup
\topic{两套贡献来源表，各自保留边界}
中美来源前十表含20条记录，2.2显示地理连接、2.6比较来源构成。扩展关系表含60条记录、14个来源与10个目的地；两表不拼接，也不以差额解释变化。
\vfill\src{所有关键例外仍在具体图文旁标注。本章用于理解结构、识别协作节点；上游身份去重、项目映射及样本定义的待核状态，详见配套核查记录。}
\clearpage
